\section{Newtonian two-dimensional bubble}\label{sec:newtonian}

\status{Measured and closed. Both methods ran the same case to $t=3$ ---
12\,000 steps, no restarts --- and the five benchmark targets are compared
against the published reference below. The central finding of the section is
not any one of the five but their \emph{sign}.}

\subsection{Case definition}

The reference case is the rising-bubble benchmark of Hysing
et~al.~\missing{cite}, which fixes the geometry, the fluid properties and the
boundary conditions and --- unusually --- publishes reference series from three
independent codes, so that the disagreement \emph{among} the references is
itself quantified. That spread is the resolution limit of the oracle: asking a
new code to do better than it would be asking for more accuracy than the
reference possesses.

\begin{adjustbox}{max width=\linewidth}
\pgfplotstabletypeset[
  string type,
  columns/quantity/.style={column name={quantity}},
  columns/TP2D/.style={column name={TP2D}},
  columns/FreeLIFE/.style={column name={FreeLIFE}},
  columns/MooNMD/.style={column name={MooNMD}},
  columns/consensus/.style={column name={consensus}},
  every head row/.style={before row=\toprule, after row=\midrule},
  every last row/.style={after row=\bottomrule},
]{\dat{hysing-referencia.dat}}
\end{adjustbox}

\subsection{The three benchmark quantities}

The benchmark is defined by three series, computed here exactly as its
Section~2.5 defines them:
%
\begin{align}
  c  &= \frac{2\sqrt{\pi A}}{P}, &
  y_c &= \frac{\int_{\text{bubble}} y \,\mathrm{d}x}{\int_{\text{bubble}} \mathrm{d}x}, &
  V_c &= \frac{\int_{\text{bubble}} u_y \,\mathrm{d}x}{\int_{\text{bubble}} \mathrm{d}x} .
\end{align}

\caveat{The circularity above is the perimeter of the area-equivalent circle
over the perimeter of the bubble. The form $4\pi A/P^2$, common in
implementations, is its \emph{square}. We used the quadratic form at one point
and compared it with the published value without converting, reporting $0.7987$
against a published $0.9013$ when the comparable number was
$\sqrt{0.7987} = 0.8937$. The discrepancy looked like a method error and was a
definition error. Both forms are called ``circularity'' in the literature.}

Both integrals are taken over the volume-fraction field --- which VOF advects
and which front-tracking writes from the geometry of the front --- so the
measurement is the same on both sides and any difference that survives belongs
to the method.

\subsection{The runs}

\begin{adjustbox}{max width=\linewidth}
\pgfplotstabletypeset[
  string type,
  columns/quantity/.style={column name={quantity}},
  columns/FT/.style={column name={front-tracking}},
  columns/VOF/.style={column name={VOF}},
  every head row/.style={before row=\toprule, after row=\midrule},
  every last row/.style={after row=\bottomrule},
]{\dat{hysing-final.dat}}
\end{adjustbox}

\subsection{The two methods agree with each other, until they do not}

Early in the run the agreement between the two methods is far closer than the
agreement of either with the reference: the centre-of-mass series differ in the
fifth decimal while both sit a few tenths of a percent below the published
values.

\begin{adjustbox}{max width=\linewidth}
\pgfplotstabletypeset[
  string type,
  columns/t/.style={column name={$t$}},
  columns/ycft/.style={column name={$y_c$ FT}},
  columns/ycvof/.style={column name={$y_c$ VOF}},
  columns/difyc/.style={column name={$\Delta y_c$}},
  columns/vcft/.style={column name={$V_c$ FT}},
  columns/vcvof/.style={column name={$V_c$ VOF}},
  columns/difvc/.style={column name={$\Delta V_c$}},
  columns/circft/.style={column name={$c$ FT}},
  every head row/.style={before row=\toprule, after row=\midrule},
  every last row/.style={after row=\bottomrule},
]{\dat{hysing-parcial.dat}}
\end{adjustbox}
\captionof{table}{Both methods over the early window.}\label{tab:parcial}


\begin{figure}[htbp]
\centering
\begin{adjustbox}{max width=\linewidth}
\begin{tikzpicture}
\begin{axis}[
  name=vc, width=7.2cm, height=5.0cm,
  xlabel={$t$}, ylabel={$V_c$}, grid=major,
  title={rise velocity, early window}, title style={font=\small},
  legend style={font=\scriptsize, at={(0.5,-0.24)}, anchor=north},
  legend columns=2, legend cell align=left,
]
\addplot[thick,mark=*,mark size=1.1]          table[x=t,y=vcft]  {\dat{hysing-parcial.dat}};
\addlegendentry{front-tracking}
\addplot[densely dashed,mark=o,mark size=1.4] table[x=t,y=vcvof] {\dat{hysing-parcial.dat}};
\addlegendentry{VOF}
\end{axis}
\begin{axis}[
  at={(vc.right of south east)}, anchor=left of south west,
  width=7.2cm, height=5.0cm,
  xlabel={$t$}, ylabel={$\Delta V_c$}, grid=major,
  ymode=log, ymin=5e-5, ymax=2e-2,
  title={divergence between the methods}, title style={font=\small},
]
\addplot[thick,mark=*,mark size=1.1] table[x=t,y=difvc] {\dat{hysing-divergencia-plot.dat}};
\end{axis}
\end{tikzpicture}
\end{adjustbox}
\caption{Left: the rise velocity of both methods over the early window, where
the two are indistinguishable --- the curves overlap to plotting accuracy and
the tabulated difference is at most \num{7.4e-5} (Table~\ref{tab:parcial}).
Right: the difference between them over the window in which it grows, on a
logarithmic axis; it grows by a factor of $68$ --- from \num{1.1e-4} to
\num{7.5e-3} --- between $t=0.4$ and $t=1.6$, which is the deformation
phase, and it is not monotone: the dip at $t=1.4$ is in the stored series. \textbf{Both panels are sampled
windows, not the full run}: the stored series cover $t \le 0.35$ and
$0.4 \le t \le 1.6$ respectively.
\missing{the continuous series over the whole run to $t=3$. It is recoverable
without new simulation --- both runs left 301 frames on disk --- but it is a
re-reduction, not a plot of stored numbers.}}
\label{fig:hysing-series}
\end{figure}

That mutual agreement does not survive the deformation. By $t=1.6$ the two
methods differ by $3.4\%$ in $V_c$, and then --- this is the part worth
noticing --- they come back together in $y_c$: $1.078587$ against $1.078645$ at
$t=3$, a difference of \num{5.8e-5}.

\begin{adjustbox}{max width=\linewidth}
\pgfplotstabletypeset[
  string type,
  columns/t/.style={column name={$t$}},
  columns/difyc/.style={column name={$\Delta y_c$}},
  columns/difvc/.style={column name={$\Delta V_c$}},
  columns/difvcrelativa/.style={column name={$\Delta V_c$ relative}},
  every head row/.style={before row=\toprule, after row=\midrule},
  every last row/.style={after row=\bottomrule},
]{\dat{hysing-divergencia.dat}}
\end{adjustbox}
\captionof{table}{Divergence between the two methods.}\label{tab:divergencia}

This is the integral-versus-derivative asymmetry: $y_c$ smooths what $V_c$
amplifies. A comparison made on centres of mass alone would conclude excellent
agreement and miss the entire discrepancy --- which is an argument for
reporting the derivative quantity even when it is the less flattering one.

\subsection{Closure: the five targets}

\begin{adjustbox}{max width=\linewidth}
\pgfplotstabletypeset[
  string type,
  columns/target/.style={column name={target}},
  columns/method/.style={column name={method}},
  columns/ours/.style={column name={ours}},
  columns/published/.style={column name={published}},
  columns/deviation/.style={column name={deviation}},
  columns/inunitsofthespread/.style={column name={in units of the spread}},
  every head row/.style={before row=\toprule, after row=\midrule},
  every last row/.style={after row=\bottomrule},
]{\dat{hysing-fechamento.dat}}
\end{adjustbox}

The last column is the one that matters: it measures the deviation in units of
the disagreement the three published codes have \emph{among themselves}.
\textbf{Three of the five fall inside it.}

\subsection{The finding that outweighs the five numbers}

\textbf{All five deviations are negative.} None on the positive side. The
bubble comes out less fast, less circular and lower than the literature, in
\emph{two} independent discretisations --- one Lagrangian and one Eulerian,
with curvature and force computations that have nothing in common.

Five out of five in the same direction is not a plausible coincidence. It is
the signature of a \textbf{shared bias}: the defect lies in what the two
methods have in common --- mesh, time step, boundary conditions, momentum
discretisation --- and not in the interface treatment.

This also corrects a conclusion we had drawn from a single target. When only
$V_{c,\max}$ had closed, the gap between front-tracking ($0.31\times$ the
spread) and VOF ($2.29\times$) pointed at the interface treatment. With all
five on the table the reading splits, and the two halves are compatible: the
\emph{difference between the methods} does answer to the interface treatment,
and in the expected direction, with the curvature of the explicit geometry
beating the curvature of a smoothed field; the \emph{common bias}, which
neither avoids, does not.

\begin{adjustbox}{max width=\linewidth}
\pgfplotstabletypeset[
  string type,
  columns/quantity/.style={column name={quantity}},
  columns/FT/.style={column name={front-tracking}},
  columns/VOF/.style={column name={VOF}},
  columns/published/.style={column name={published}},
  every head row/.style={before row=\toprule, after row=\midrule},
  every last row/.style={after row=\bottomrule},
]{\dat{hysing-vcmax.dat}}
\end{adjustbox}

Three candidates for the bias, \textbf{none of them measured}: the single
resolution (40 cells per bubble radius), the capillary time-step restriction,
and the possibility that the wake refinement is insufficient. A convergence
study in $h$ separates the three, and it is the first item on the list of next
steps --- now with a measured reason to be first rather than for methodological
completeness.

\subsection{Shape}
\begin{figure}[htbp]
\centering
\begin{tikzpicture}
\begin{axis}[
  name=evol, width=6.2cm, height=9.2cm, axis equal image,
  xlabel={$x$}, ylabel={$y$}, grid=major,
  xmin=0.11, xmax=0.89, ymin=0.22, ymax=1.32,
  title={the bubble, $t=0$ to $t=3$}, title style={font=\small},
  legend style={font=\scriptsize, at={(0.5,-0.13)}, anchor=north},
  legend columns=2, legend cell align=left,
]
\addplot[thick]           table[x=x,y=y] {\dat{formas/hy_ft_0.dat}};
\addlegendentry{front-tracking}
\addplot[densely dashed]  table[x=x,y=y] {\dat{formas/hy_vof_0.dat}};
\addlegendentry{VOF}
\addplot[thick]           table[x=x,y=y] {\dat{formas/hy_ft_1.dat}};
\addplot[densely dashed]  table[x=x,y=y] {\dat{formas/hy_vof_1.dat}};
\addplot[thick]           table[x=x,y=y] {\dat{formas/hy_ft_2.dat}};
\addplot[densely dashed]  table[x=x,y=y] {\dat{formas/hy_vof_2.dat}};
\addplot[thick]           table[x=x,y=y] {\dat{formas/hy_ft_3.dat}};
\addplot[densely dashed]  table[x=x,y=y] {\dat{formas/hy_vof_3.dat}};
\end{axis}
\begin{axis}[
  at={(evol.right of south east)}, anchor=left of south west,
  width=6.6cm, height=6.6cm, axis equal image,
  xlabel={$x$}, grid=major,
  xmin=0.147, xmax=0.167, ymin=0.794, ymax=0.814,
  title={$t=1.9$, magnified 34 times}, title style={font=\small},
  scaled ticks=false, tick label style={font=\tiny, /pgf/number format/fixed,
                                        /pgf/number format/precision=3},
  xtick={0.150,0.155,0.160,0.165}, ytick={0.795,0.800,0.805,0.810},
]
\addplot[thick]           table[x=x,y=y] {\dat{formas/hy_ft_2.dat}};
\addplot[densely dashed]  table[x=x,y=y] {\dat{formas/hy_vof_2.dat}};
\end{axis}
\end{tikzpicture}
\caption{The interface in the two methods at $t=0$, $1$, $1.9$ and $3$. The
bubble rises from $y_c=0.5$ to $1.08$ and flattens: the width-to-height ratio
goes from $1.00$ to $1.87$, and $t=1.9$ is the instant of maximum deformation
--- the same instant as the minimum circularity. From $1.9$ to $3$ the shape
barely changes ($1.87$ to $1.86$): the bubble reaches a steady regime. The two
curves are \emph{indistinguishable} over the whole evolution; the right-hand
panel magnifies 34 times the point of \emph{greatest} discrepancy at the
instant of maximum deformation, where they separate by \num{4.7e-4} --- about
\SI{8}{\percent} of one fine cell. The VOF curve is the
$\mathrm{FracVol}=0.5$ contour by marching squares.}
\label{fig:hysing-formas}
\end{figure}


\begin{adjustbox}{max width=\linewidth}
\pgfplotstabletypeset[
  string type,
  columns/t/.style={column name={$t$}},
  columns/razaolarguraaltura/.style={column name={width/height}},
  columns/areaft/.style={column name={$A$ FT}},
  columns/areavof/.style={column name={$A$ VOF}},
  columns/perimft/.style={column name={$P$ FT}},
  columns/perimvof/.style={column name={$P$ VOF}},
  columns/cdosdois/.style={column name={$c$, both}},
  every head row/.style={before row=\toprule, after row=\midrule},
  every last row/.style={after row=\bottomrule},
]{\dat{hysing-forma.dat}}
\end{adjustbox}

\subsection{Under an adaptive mesh, VOF loses exact conservation}

On a fixed mesh the VOF fraction is conserved to machine precision. Under
adaptive refinement it is not, and the loss is not noise: the samples are
almost monotone.

\begin{adjustbox}{max width=\linewidth}
\pgfplotstabletypeset[
  string type,
  columns/method/.style={column name={method}},
  columns/fixed-mesh/.style={column name={fixed mesh}},
  columns/with-AMR-until-t/.style={column name={with AMR, until $t$}},
  columns/monotone-samples/.style={column name={monotone samples}},
  every head row/.style={before row=\toprule, after row=\midrule},
  every last row/.style={after row=\bottomrule},
]{\dat{hysing-massa-amr.dat}}
\end{adjustbox}

Front-tracking, which does not store the fraction, is affected far less. The
asymmetry is a property of the representation and not of the refinement
machinery, which is shared.

\subsection{Case 2: the skirted regime}

\missing{not run. Worth doing precisely because the published reference
solutions disagree with each other there: it is the regime that discriminates
between interface representations rather than confirming both.}
